Um die Gleichheit $\mathrm{GL}_n(\mathbb{C}) \cap \mathrm{O}(2n) = \mathrm{U}(n)$ zu zeigen, betrachten wir die Definitionen der beteiligten Gruppen. Die allgemeine lineare Gruppe $\mathrm{GL}_n(\mathbb{C})$ besteht aus allen invertierbaren $n \times n$ Matrizen mit komplexen Einträgen. Die orthogonale Gruppe $\mathrm{O}(2n)$ besteht aus allen $2n \times 2n$ Matrizen $A$ mit reellen Einträgen, die die Bedingung $A^T A = I_{2n}$ erfüllen, wobei $I_{2n}$ die $2n \times 2n$ Einheitsmatrix ist. Die unitäre Gruppe $\mathrm{U}(n)$ besteht aus allen $n \times n$ Matrizen $U$ mit komplexen Einträgen, die die Bedingung $U^* U = I_n$ erfüllen, wobei $U^*$ die konjugiert-transponierte Matrix von $U$ ist.

Um zu zeigen, dass $\mathrm{GL}_n(\mathbb{C}) \cap \mathrm{O}(2n) = \mathrm{U}(n)$, müssen wir zeigen, dass eine Matrix $A \in \mathrm{GL}_n(\mathbb{C})$ genau dann in $\mathrm{O}(2n)$ liegt, wenn sie in $\mathrm{U}(n)$ liegt. Eine Matrix $A \in \mathrm{GL}_n(\mathbb{C})$ liegt in $\mathrm{O}(2n)$, wenn $A^T A = I_{2n}$. Da $A$ komplex ist, bedeutet dies, dass $A^* A = I_n$, was genau die Bedingung für $A \in \mathrm{U}(n)$ ist. Daher ist $\mathrm{GL}_n(\mathbb{C}) \cap \mathrm{O}(2n) = \mathrm{U}(n)$.

Für die zweite Gleichheit $\mathrm{GL}_n(\mathbb{R}) \cap \mathrm{U}(n) = \mathrm{O}(n)$ betrachten wir die Definitionen der Gruppen. Die Gruppe $\mathrm{GL}_n(\mathbb{R})$ besteht aus allen invertierbaren $n \times n$ Matrizen mit reellen Einträgen. Eine Matrix $A \in \mathrm{GL}_n(\mathbb{R})$ liegt in $\mathrm{U}(n)$, wenn $A^* A = I_n$. Da $A$ reell ist, ist $A^* = A^T$, sodass die Bedingung $A^* A = I_n$ äquivalent zu $A^T A = I_n$ ist, was genau die Bedingung für $A \in \mathrm{O}(n)$ ist. Daher ist $\mathrm{GL}_n(\mathbb{R}) \cap \mathrm{U}(n) = \mathrm{O}(n)$.